% !TeX program = pdflatex
% SEDAC P2.02 MC01 - Guia explicada en LaTeX Beamer
% Fuente unica. Compatible con pdfLaTeX y Tectonic.
\documentclass[10pt,aspectratio=169]{beamer}
\usepackage[T1]{fontenc}
\usepackage[utf8]{inputenc}
\usepackage[spanish,es-nodecimaldot]{babel}
\usepackage{lmodern}
\usepackage{amsmath,amssymb,booktabs,array}
\usefonttheme{professionalfonts}
\usefonttheme[onlymath]{serif}
\definecolor{azul}{HTML}{173345}
\definecolor{verde}{HTML}{007D86}
\definecolor{gris}{HTML}{576B79}
\definecolor{rojo}{HTML}{9D302C}
\setbeamercolor{normal text}{fg=azul,bg=white}
\setbeamercolor{structure}{fg=verde}
\setbeamercolor{frametitle}{fg=azul}
\setbeamercolor{title}{fg=azul}
\setbeamercolor{alerted text}{fg=rojo}
\setbeamerfont{frametitle}{size=\Large,series=\bfseries}
\setbeamerfont{title}{size=\LARGE,series=\bfseries}
\setbeamertemplate{navigation symbols}{}
\setbeamertemplate{itemize item}{\small\textbullet}
\setbeamersize{text margin left=9mm,text margin right=9mm}
\setlength{\parskip}{4pt}
\newcommand{\mc}[1]{\texttt{\detokenize{#1}}}
\newcommand{\N}{\,\mathrm{N}}
\newcommand{\mm}{\,\mathrm{mm}}
\newcommand{\MPa}{\,\mathrm{MPa}}
\newcommand{\kg}{\,\mathrm{kg}}
\newcommand{\cita}[1]{\par\medskip{\scriptsize\color{gris}#1}}
\newcommand{\mcref}{Guía de lectura}
\setbeamertemplate{footline}{%
 \leavevmode\hbox{%
 \begin{beamercolorbox}[wd=.37\paperwidth,ht=2.8ex,dp=1.25ex,leftskip=9mm]{normal text}
 \scriptsize\color{gris}SEDAC P2.02 / MC01
 \end{beamercolorbox}%
 \begin{beamercolorbox}[wd=.46\paperwidth,ht=2.8ex,dp=1.25ex,center]{normal text}
 \scriptsize\color{verde}\mcref
 \end{beamercolorbox}%
 \begin{beamercolorbox}[wd=.17\paperwidth,ht=2.8ex,dp=1.25ex,rightskip=9mm]{normal text}
 \hfill\scriptsize\color{gris}\insertframenumber/\inserttotalframenumber
 \end{beamercolorbox}}}
\title{Memoria de cálculo explicada}
\subtitle{SEDAC P2.02\\Desarrollos y criterios de falla}
\author{Guía para acompañar la hoja MC01 de Mathcad}
\date{15 de octubre de 2026}

\begin{document}
\spanishplainpercent
\gdef\mcref{Guía de lectura}
\begin{frame}[label=inicio]
 \titlepage
 \vspace{-5mm}
 \begin{center}
  \small Prediseño paramétrico sujeto a mediciones de la bomba real.
 \end{center}
\end{frame}

\begin{frame}[t,label=mapa]{Índice para seguir Mathcad}
\small
Cada número corresponde a la \textbf{sección de la hoja MC01}.
Los enlaces llevan al primer desarrollo de ese tema.
\medskip
\begin{columns}[T,onlytextwidth]
\column{.49\textwidth}
\begin{enumerate}
 \item \hyperlink{mc01}{Alcance y cargas}
 \item \hyperlink{mc02}{Geometría y brazos}
 \item \hyperlink{mc03}{Materiales y factores}
 \item \hyperlink{mc04}{Sección redondeada}
 \item \hyperlink{mc05}{Riel: axial y flexión}
 \item \hyperlink{mc06}{Riel: torsión, fluencia y flecha}
 \item \hyperlink{mc07}{Estabilidad}
 \item \hyperlink{mc08}{Grupo de tornillos M8}
 \item \hyperlink{mc09}{Cortante y falla de chapa}
\end{enumerate}
\column{.49\textwidth}
\begin{enumerate}
 \setcounter{enumi}{9}
 \item \hyperlink{mc10}{Raíz de cartela}
 \item \hyperlink{mc11}{Cargas sobre el piso}
 \item \hyperlink{mc12}{Piso: tensión y espesor}
 \item \hyperlink{mc13}{Refuerzos y hombro}
 \item \hyperlink{mc14}{Cinturón y diagonales}
 \item \hyperlink{mc15}{M6 y alcance de los criterios}
 \item \hyperlink{mc16}{Resumen de resultados}
 \item \hyperlink{mc17}{Actualización del modelo}
\end{enumerate}
\end{columns}
\medskip
\footnotesize
Antes de los cálculos: operadores, funciones y fluencia.
Al final: fatiga, fractura, sensibilidad, diccionario y fuentes.
\end{frame}

\begin{frame}[t]{Notación y uso de esta guía}
\begin{columns}[T,onlytextwidth]
\column{.52\textwidth}
\textbf{En las ecuaciones}
\[
 \sigma_{\mathrm{VM}},\qquad F_v,\qquad
 n_R=\frac{S_{yR}}{\sigma_{\mathrm{VM}}}
\]
La notación matemática usa letras griegas y subíndices para facilitar la lectura.
\medskip

\textbf{En Mathcad}
\[
 \mc{SigVM},\qquad\mc{Fv},\qquad\mc{nR}
\]
Debajo de cada desarrollo se indican los nombres que aparecen en la hoja.
\column{.44\textwidth}
\textbf{Cómo seguir los desarrollos}
\begin{itemize}
 \item Abrir la hoja editable MC01.
 \item Buscar el mismo número de sección en ambos archivos.
 \item Comparar la expresión general con la sustitución numérica.
 \item Revisar la hipótesis antes de cambiar un dato.
\end{itemize}
\end{columns}
\medskip
\textbf{Los valores del PDF son estáticos.}
Al editar Mathcad, sus resultados pueden diferir de los ejemplos de esta guía.
\end{frame}

\begin{frame}[t]{Mathcad: definición y evaluación}
\begin{columns}[T,onlytextwidth]
\column{.48\textwidth}
\textbf{Definición con \(\boldsymbol{:=}\)}
\[
 \mc{Mass}:=25\kg
\]
\[
 \mc{Fv}:=\mc{Mass}\,\mc{grav}\,\mc{av}
\]
Asigna un valor o una regla. Una definición puede estar bien construida sin mostrar un resultado al lado.
\column{.48\textwidth}
\textbf{Evaluación con \(\boldsymbol{=}\)}
\[
 \mc{Fv}=735{,}75\N
\]
Muestra el valor de una expresión ya definida.
\medskip

\textbf{Orden de cálculo}

De arriba hacia abajo y, a igual altura, de izquierda a derecha.
\end{columns}
\bigskip
\textbf{Recálculo completo:} la ayuda de Mathcad clásico indica
\mc{Ctrl + F9}. En modo manual puede aparecer \mc{Calc} en la barra de estado.
\cita{[1] Ayuda instalada de Mathcad clásico: operadores, orden de cálculo y cálculo automático.}
\end{frame}

\begin{frame}[t]{Mathcad: funciones y unidades}
\begin{columns}[T,onlytextwidth]
\column{.52\textwidth}
\textbf{Funciones utilizadas}
\[
 \sqrt{9\mm^2}=3\mm
\]
\[
 \min(I_x,I_z)\quad\text{elige el eje menos rígido}
\]
\[
 \max(M_{\mathrm{tramo}},M_{\mathrm{borde}})
\]
La última expresión selecciona el caso de piso más exigente.

\textbf{Mayúsculas y minúsculas}
\[
 \mc{Nr}\ [\mathrm{N}]\neq\mc{nR}\ [-]
\]
Una es fuerza axial y la otra es un margen.
\column{.44\textwidth}
\textbf{Unidades y potencias}
\[
 1\MPa=1\,\frac{\mathrm{N}}{\mathrm{mm}^2}
\]
\[
 1\mm^4=10^{-12}\,\mathrm{m}^4
\]
\[
 I_x=3{,}2784\times10^{-8}\,\mathrm{m}^4
      =32\,784{,}01\mm^4
\]
Introducir una longitud requiere una unidad:
\(\mc{t}:=2{,}5\mm\).
\end{columns}
\cita{[1] La unidad de presentación cambia la forma del resultado, manteniendo su valor físico.}
\end{frame}

\begin{frame}[t]{Fluencia: Von Mises y margen de cálculo}
\textbf{Hipótesis de material dúctil.}
Se compara una tensión equivalente con la referencia de fluencia \(S_y\).
\[
 \boxed{\sigma_{\mathrm{VM}}=\sqrt{\sigma^2+3\tau^2}}
 \qquad
 \boxed{n=\frac{S_y}{\sigma_{\mathrm{VM}}}}
\]
La primera expresión corresponde a un estado con una tensión normal
\(\sigma\) y una cortante \(\tau\) en el mismo punto.
\begin{columns}[T,onlytextwidth]
\column{.48\textwidth}
\textbf{Cortante puro}
\[
 \sigma=0
 \quad\Longrightarrow\quad
 \sigma_{\mathrm{VM}}=\sqrt{3}\,|\tau|
\]
De aquí proviene \(\sqrt{3}\) en las comprobaciones M8 y M6.
\column{.48\textwidth}
\textbf{Objetivo adoptado}
\[
 n\ge n_{\mathrm{req}}=2{,}2
\]
Este valor es una decisión provisional del prediseño.
\end{columns}
\cita{[2] MIT, Structural Plasticity. El riel usa una envolvente de solicitaciones que debe revisarse localmente.}
\end{frame}

\begin{frame}[t]{Criterios de falla y factores del modelo}
\begin{columns}[T,onlytextwidth]
\column{.52\textwidth}
\textbf{Tresca, como referencia teórica}
\[
 \sigma_{\mathrm{T}}=
 \max\!\left(
 |\sigma_1-\sigma_2|,
 |\sigma_2-\sigma_3|,
 |\sigma_3-\sigma_1|
 \right)
\]
En cortante puro:
\[
 \sigma_{\mathrm{T}}=2|\tau|,
 \qquad
 \sigma_{\mathrm{VM}}=\sqrt{3}\,|\tau|
\]
La hoja utiliza Von Mises. Tresca se explica aquí para distinguir ambos criterios.
\column{.44\textwidth}
\textbf{Los números cumplen funciones distintas}
\begin{itemize}
 \item \(a_v,a_h\): niveles de carga elegidos.
 \item \(K_t,K_{\mathrm{slot}},K_b\): amplificaciones provisionales de demanda.
 \item \(n_{\mathrm{req}}\): margen buscado.
\end{itemize}
Un margen \(n=1{,}22\) queda por debajo del objetivo \(2{,}2\), aunque sea mayor que uno.
\end{columns}
\cita{[2] Fluencia, pandeo, fatiga, fractura y deformación excesiva requieren comprobaciones diferentes.}
\end{frame}

\section{1. Alcance y cargas}
\gdef\mcref{Mathcad / sección 1}
\begin{frame}[t,label=mc01]{1. Fuerzas de cálculo}
\begin{columns}[T,onlytextwidth]
\column{.49\textwidth}
\textbf{Entradas provisionales}
\[
 m=25\kg,\qquad g=9{,}81\,\mathrm{m/s^2}
\]
\[
 a_v=3,\qquad a_h=1,\qquad\eta=0{,}7
\]
\(a_v=3\) representa tres veces el peso, incluyendo la gravedad.
\column{.47\textwidth}
\textbf{Sustitución}
\begin{align*}
 F_v&=mga_v\\
    &=25(9{,}81)(3)=735{,}75\N\\
 F_h&=mga_h=245{,}25\N\\
 N_R&=\eta F_v=515{,}025\N
\end{align*}
\end{columns}
\medskip
\(\eta\) representa el lado que recibe el \(70\,\%\) de la carga en un
reparto supuesto \(70/30\). Debe revisarse con el arnés y la bomba reales.
\cita{Mathcad: \mc{Mass}, \mc{grav}, \mc{av}, \mc{ah}, \mc{eta}, \mc{Fv}, \mc{Fh}, \mc{Nr}.
La masa de 25 kg es una hipótesis de cálculo.}
\end{frame}

\section{2. Geometría}
\gdef\mcref{Mathcad / sección 2}
\begin{frame}[t,label=mc02]{2. Geometría y brazos de palanca}
\begin{columns}[T,onlytextwidth]
\column{.48\textwidth}
\textbf{Riel tubular}
\[
 b=50\mm,\quad d=25\mm,\quad t=2{,}5\mm
\]
\[
 H=500\mm,\qquad s=200\mm
\]
\[
 r_o=5\mm,\qquad r_i=2{,}5\mm
\]
\textbf{Excentricidad}
\[
 e=100\mm
\]
Se mide desde el eje del riel hasta la línea de acción de la carga vertical.
\column{.48\textwidth}
\textbf{Longitud idealizada}
\begin{align*}
 L&=\max(L_{\mathrm{input}},H-100\mm)\\
  &=\max(400,400)\mm\\
  &=400\mm
\end{align*}
\[
 n_{\mathrm{req}}=2{,}2,\qquad
 \delta_{\mathrm{lim}}=5\mm
\]
La regla de \(L\) debe contrastarse con los apoyos reales.
\end{columns}
\cita{Mathcad: \mc{b}, \mc{d}, \mc{t}, \mc{H}, \mc{Sep}, \mc{ecc}, \mc{Le}.
La geometría de cálculo corresponde a escala 1:1.}
\end{frame}

\section{3. Materiales}
\gdef\mcref{Mathcad / sección 3}
\begin{frame}[t,label=mc03]{3. Propiedades y factores adoptados}
\begin{columns}[T,onlytextwidth]
\column{.56\textwidth}
\begin{tabular}{@{}lll@{}}
\toprule
Propiedad & Valor & Uso\\
\midrule
\(S_{yR}\) & \(240\MPa\) & Perfil 6061-T6\\
\(S_{yS}\) & \(130\MPa\) & Chapa 5052-H32\\
\(S_{yB}\) & \(450\MPa\) & Tornillo A2-70\\
\(E\) & \(69\,000\MPa\) & Rigidez\\
\(\nu\) & \(0{,}33\) & Poisson\\
\(\rho\) & \(2\,700\,\mathrm{kg/m^3}\) & Masa\\
\bottomrule
\end{tabular}
\[
 K_t=2,\quad K_{\mathrm{slot}}=1{,}5,\quad K_b=1{,}2
\]
\column{.40\textwidth}
\(S_y\) mide resistencia a fluencia.

\(E\) relaciona tensión con deformación elástica.

El producto, el espesor y el temple deben coincidir con el certificado del material.

Los factores \(K\) son provisionales. Aún faltan modelos locales para obtener factores específicos de cada detalle.
\end{columns}
\cita{[3--6] Hydro, Atlas y Bossard. Para 5052-H32 se adopta 130 MPa considerando diferencias entre especificaciones.
Una zona soldada requiere revisar la resistencia del 6061-T6.}
\end{frame}

\section{4. Sección redondeada}
\gdef\mcref{Mathcad / sección 4}
\begin{frame}[t,label=mc04]{4. Área del tubo con radios de esquina}
Se retiran cuatro cuadrados y se restituyen cuatro cuartos de círculo:
\[
 \boxed{\operatorname{Ar}(B,D,R)=BD-(4-\pi)R^2}
\]
El tubo resulta de restar el hueco interior:
\[
 A=\operatorname{Ar}(b,d,r_o)
 -\operatorname{Ar}(b-2t,d-2t,r_i)
\]
\[
 \boxed{A=333{,}90486\mm^2}
\]
\textbf{Función frente a resultado}
\[
 \operatorname{Ar}(50\mm,25\mm,5\mm)=1\,228{,}54\mm^2
\]
Esta última cifra corresponde al rectángulo redondeado macizo.
\(A\) es el área de metal del tubo.
\cita{Mathcad: función \mc{Ar(B,D,R)} y variable \mc{Area}.
Los argumentos \(B,D,R\) son locales a la función.}
\end{frame}

\begin{frame}[t]{4. Desarrollo del segundo momento de esquina}
En una esquina se integra \((c+v)^2\) sobre el área retirada, con
\(c=D/2-R\).
\[
 c_0=1-\frac{\pi}{4},
 \qquad
 c_1=\frac13-\frac{\pi}{16}
\]
\[
 \int_{\mathrm{ret}}dA=c_0r^2,\qquad
 \int_{\mathrm{ret}}v\,dA=\frac{r^3}{6},\qquad
 \int_{\mathrm{ret}}v^2\,dA=c_1r^4
\]
Al expandir el cuadrado:
\begin{align*}
 J_c(c,r)
 &=c^2\!\int_{\mathrm{ret}}dA
   +2c\!\int_{\mathrm{ret}}v\,dA
   +\int_{\mathrm{ret}}v^2\,dA\\
 &=\boxed{c_0c^2r^2+\frac{cr^3}{3}+c_1r^4}
\end{align*}
\cita{Mathcad: \mc{Jc(c,r)}. Es un segundo momento de área auxiliar, de unidad longitud\(^{4}\).
Su nombre no corresponde a la constante de torsión de Saint-Venant.}
\end{frame}

\begin{frame}[t]{4. Inercias del rectángulo y del hueco}
Para el eje paralelo a \(B\):
\[
 \operatorname{Ir}(B,D,R)
 =\frac{BD^3}{12}
 -4J_c\!\left(\frac D2-R,R\right)
\]
\begin{align*}
 I_x&=\operatorname{Ir}(b,d,r_o)
 -\operatorname{Ir}(b-2t,d-2t,r_i),\\
 I_z&=\operatorname{Ir}(d,b,r_o)
 -\operatorname{Ir}(d-2t,b-2t,r_i).
\end{align*}
\[
 \boxed{I_x=32\,784{,}01\mm^4}
 \qquad
 \boxed{I_z=98\,865{,}89\mm^4}
\]
Intercambiar \(B\) y \(D\) cambia el eje de flexión.

Estas propiedades describen la sección continua idealizada.
Los agujeros y uniones requieren revisar la sección local.
\cita{Mathcad: \mc{Ir(B,D,R)}, \mc{Ix}, \mc{Iz}.}
\end{frame}

\section{5. Riel: tensión normal}
\gdef\mcref{Mathcad / sección 5}
\begin{frame}[t,label=mc05]{5. Esfuerzo axial y flexión por excentricidad}
\begin{columns}[T,onlytextwidth]
\column{.48\textwidth}
\textbf{Módulos resistentes}
\[
 Z_x=\frac{I_x}{d/2},\qquad
 Z_z=\frac{I_z}{b/2}
\]
De \(\sigma=My/I\), en la fibra exterior:
\[
 \sigma_{\max}=\frac{Mc}{I}=\frac MZ
\]
\textbf{Carga del lado crítico}
\[
 N_R=\eta F_v,\qquad M_e=\eta F_ve
\]
\column{.48\textwidth}
\textbf{Resultados}
\begin{align*}
 N_R&=515{,}025\N\\
 M_e&=51\,502{,}50\N\mm\\
 \sigma_{\mathrm{ax}}&=\frac{N_R}{A}
                       =1{,}54\MPa\\
 \sigma_e&=\frac{M_e}{Z_x}
                       =19{,}64\MPa
\end{align*}
\end{columns}
\[
 m_{\mathrm{rieles}}=2HA\rho=0{,}902\kg
\]
La masa anterior corresponde a los tubos idealizados sin descontar agujeros.
\cita{Mathcad: \mc{Zx}, \mc{Zz}, \mc{Nr}, \mc{Me}, \mc{SigAx}, \mc{SigEc}, \mc{MassR}.}
\end{frame}

\begin{frame}[t]{5. Máximo de flexión por fuerza horizontal}
Se permite una dirección horizontal cualquiera de módulo \(F_h\).
Para dos componentes perpendiculares \(F_1,F_2\):
\[
 F_1^2+F_2^2=F_h^2,
 \qquad
 \sigma_h=\eta L\!\left(\frac{F_1}{Z_x}+\frac{F_2}{Z_z}\right).
\]
La desigualdad de Cauchy-Schwarz da:
\[
 \boxed{\sigma_h=
 \eta F_hL\sqrt{\frac1{Z_x^2}+\frac1{Z_z^2}}}
 =31{,}42\MPa
\]
Se suman las contribuciones que agravan la fibra crítica:
\[
 \sigma_N=\sigma_{\mathrm{ax}}+\sigma_e+\sigma_h
          =\boxed{52{,}60\MPa}
\]
El resultado corresponde al modelo de riel en voladizo.
Las reacciones del arnés pueden cambiar esta distribución.
\cita{Mathcad: \mc{SigH}, \mc{SigN}.}
\end{frame}

\section{6. Riel: torsión y flecha}
\gdef\mcref{Mathcad / sección 6}
\begin{frame}[t,label=mc06]{6. Torsión del tubo cerrado}
\textbf{Área encerrada por la línea media de la pared}
\[
 A_m=(b-t)(d-t)-(4-\pi)(1{,}5t)^2
     =1\,056{,}68\mm^2
\]
Para una celda delgada cerrada, el flujo cortante es aproximadamente
\(q_T=T/(2A_m)\). Al dividir por el espesor:
\[
 T_R=\eta F_h\frac{s}{2}=17\,167{,}50\N\mm
\]
\[
 \boxed{\tau_T=\frac{T_R}{2A_mt}=3{,}25\MPa}
\]
\(A_m\) encierra la línea media del tubo.
\(A=333{,}90\mm^2\) representa el metal de su pared.

La aproximación supone una celda continua y espesor uniforme.
Las perforaciones pueden alterar el flujo cortante.
\cita{Mathcad: \mc{Am}, \mc{Tr}, \mc{TauT}.}
\end{frame}

\begin{frame}[t]{6. Envolvente de tensiones y fluencia del riel}
La hoja conserva la siguiente aproximación de cortante medio:
\[
 \tau_V=\frac{\eta\sqrt{F_v^2+F_h^2}}{A}
       =1{,}63\MPa
\]
Con \(K_t=2\):
\[
 \sigma_{\mathrm{VM}}
 =K_t\sqrt{\sigma_N^2+3(\tau_T+\tau_V)^2}
 =106{,}54\MPa
\]
\[
 \boxed{n_R=\frac{240}{106{,}54}=2{,}253}
 \qquad n_{\mathrm{req}}=2{,}2
\]
\footnotesize
\(\tau_V\) incluye \(F_v\), aunque el riel ya la recibe axialmente.
\(V/A\) es un valor medio. La suma combina máximos de una envolvente
y necesita revisión como campo local, especialmente cerca de agujeros.

El margen supera el objetivo en aproximadamente \(2{,}4\,\%\).
\cita{Mathcad: \mc{TauV}, \mc{SigVM}, \mc{nR}.}
\end{frame}

\begin{frame}[t]{6. Flecha del riel idealizado}
Para una viga elástica de rigidez constante:
\[
 EI\,v''(x)=M(x).
\]
Integrando con desplazamiento y giro nulos en el empotramiento:
\begin{align*}
 \delta_h&=\frac{\eta F_hL^3}{3E\min(I_x,I_z)}
           =1{,}619\mm,\\
 \delta_e&=\frac{M_eL^2}{2EI_x}
           =1{,}821\mm.
\end{align*}
\[
 \boxed{\delta_R=\delta_h+\delta_e=3{,}440\mm<5\mm}
\]
La suma es una aproximación de desplazamiento.
Aún falta la flexibilidad de uniones y arnés.
\cita{[7] MIT, Beam Deflections.
Mathcad: \mc{deltaH}, \mc{deltaE}, \mc{deltaR}, \mc{deltaLim}.}
\end{frame}

\section{7. Estabilidad}
\gdef\mcref{Mathcad / sección 7}
\begin{frame}[t,label=mc07]{7. Pandeo global de Euler}
\[
 \boxed{P_{\mathrm{cr}}=
 \frac{\pi^2E I_{\min}}{(K_{\mathrm{col}}L)^2}},
 \qquad I_{\min}=\min(I_x,I_z)
\]
\[
 K_{\mathrm{col}}=2,\qquad L=400\mm
\]
\[
 P_{\mathrm{cr}}=34\,884{,}37\N,
 \qquad
 n_{\mathrm{Euler}}=\frac{P_{\mathrm{cr}}}{N_R}=67{,}73
\]
\textbf{Condiciones del modelo}
\begin{itemize}
 \item Columna ideal, recta y elástica.
 \item Un extremo empotrado y otro libre para \(K_{\mathrm{col}}=2\).
 \item Compresión y eje débil según las hipótesis adoptadas.
\end{itemize}
Ese cociente grande convive con otros modos que pueden limitar antes.
\cita{[8] MIT, Buckling of Beams.
Mathcad: \mc{Kcol}, \mc{Pcr}, \mc{nEuler}.}
\end{frame}

\begin{frame}[t]{7. Pandeo local de una pared}
Se aproxima una pared por una placa larga apoyada en sus bordes.
\[
 b_p=\max(b,d)-4t=40\mm,\qquad k=4
\]
\[
 \boxed{\sigma_{\mathrm{cr}}=
 \frac{k\pi^2E}{12(1-\nu^2)}
 \left(\frac{t}{b_p}\right)^2}
 =995{,}09\MPa
\]
Como la predicción elástica supera la fluencia adoptada:
\[
 \sigma_{\mathrm{cap}}=
 \min(S_{yR},\sigma_{\mathrm{cr}})
 =240\MPa
\]
El coeficiente \(k\), el ancho plano, los agujeros y las condiciones
de apoyo necesitan revisión en la sección real.

Quedan pendientes imperfecciones y modos de estabilidad de las uniones.
\cita{[9] MIT, Buckling of Plates and Sections.
Mathcad: \mc{bp}, \mc{kplate}, \mc{SigCr}, \mc{SigCap}.}
\end{frame}

\section{8. Grupo M8}
\gdef\mcref{Mathcad / sección 8}
\begin{frame}[t,label=mc08]{8. Reparto en dos tornillos M8}
\begin{columns}[T,onlytextwidth]
\column{.47\textwidth}
\textbf{Fuerzas directas}
\[
 V_y=\eta F_v,\qquad V_z=\eta F_h
\]
Cada tornillo recibe \(V_y/2\) y \(V_z/2\).

\textbf{Momento}
\[
 M_c=\eta(F_ve+F_ha_H)
\]
\[
 p=45\mm,\qquad a_H=62{,}5\mm
\]
\column{.49\textwidth}
\textbf{Par de fuerzas}
\[
 Qp=M_c
 \quad\Longrightarrow\quad Q=\frac{M_c}{p}
\]
En el tornillo crítico se suman las componentes del mismo sentido:
\[
 \boxed{
 R_8=K_b\sqrt{
 \left(\frac{V_y}{2}\right)^2+
 \left(\frac{V_z}{2}+\frac{M_c}{p}\right)^2}}
\]
\end{columns}
\medskip
Se idealiza un grupo de igual rigidez en el plano \(YZ\).
El otro tornillo recibe la contribución del par con signo opuesto.
\cita{Mathcad: \mc{Vy}, \mc{Vz}, \mc{MomC}, \mc{pitch}, \mc{armH}, \mc{R8}.}
\end{frame}

\begin{frame}[t]{8. Sustitución y corrección del M8}
\begin{align*}
 \frac{V_y}{2}&=257{,}5125\N,&
 \frac{V_z}{2}&=85{,}8375\N,\\
 M_c&=62\,232{,}19\N\mm,&
 \frac{M_c}{p}&=1\,382{,}94\N.
\end{align*}
\[
 R_8=1{,}2\sqrt{257{,}5125^2+
                (85{,}8375+1\,382{,}94)^2}
    =\boxed{1\,789{,}41\N}
\]
La comparación anterior de la hoja conserva:
\[
 R_{\mathrm{old}}=
 K_b\sqrt{
 \left(\frac{\sqrt{V_y^2+V_z^2}}{2}\right)^2+
 \left(\frac{M_c}{p}\right)^2}
 =1\,691{,}19\N
\]
La resistencia se revisa con \(R_8\).
Quedan por estudiar tracción, palanca y apertura fuera del plano.
\cita{[10] NASA RP-1228 como referencia de diseño de uniones.
Mathcad: \mc{R8old}, \mc{R8}.}
\end{frame}

\section{9. M8 y falla de chapa}
\gdef\mcref{Mathcad / sección 9}
\begin{frame}[t,label=mc09]{9. Cortante del tornillo y aplastamiento}
\begin{columns}[T,onlytextwidth]
\column{.47\textwidth}
\textbf{Cortante del M8}
\[
 A_{s8}=36{,}6\mm^2
\]
\[
 \sigma_{\mathrm{VM8}}
 =\sqrt{3}\,\frac{R_8}{A_{s8}}
 =84{,}68\MPa
\]
\[
 \boxed{n_{\mathrm{M8}}=\frac{450}{84{,}68}=5{,}31}
\]
Se usa área resistente de rosca y un plano de corte.
\column{.49\textwidth}
\textbf{Aplastamiento medio}
\[
 \sigma_{\mathrm{bear}}=\frac{R_8}{d_8t_{\mathrm{pieza}}}
\]
\[
 \begin{array}{c|c|c}
 &\sigma_{\mathrm{bear}}&n\\ \hline
 \text{Tubo, }t=2{,}5\mm&89{,}47\MPa&2{,}68\\
 \text{Cartela, }t_g=4\mm&55{,}92\MPa&2{,}32
 \end{array}
\]
La comprobación acredita una pared del tubo.
\end{columns}
\medskip
Comparar presión media con \(S_y\) es un precriterio.
Una unión completa requiere valores y geometrías de contacto apropiados.
\cita{Mathcad: \mc{VM8}, \mc{nBolt8}, \mc{BearR}, \mc{nBearR}, \mc{BearG}, \mc{nBearG}.}
\end{frame}

\begin{frame}[t]{9. Arrancamiento hacia un borde}
Se consideran dos planos de cortante en la chapa:
\[
 e_b=\min(15\mm,d/2)=12{,}5\mm,\qquad d_h=8{,}6\mm
\]
\[
 A_{\mathrm{arr}}=
 2\left(e_b-\frac{d_h}{2}\right)t_g
 =2(12{,}5-4{,}3)(4)=65{,}6\mm^2
\]
\[
 \tau_{\mathrm{arr}}=\frac{R_8}{A_{\mathrm{arr}}},
 \qquad
 \sigma_{\mathrm{eq,arr}}=\sqrt{3}\,\tau_{\mathrm{arr}}
\]
\[
 \boxed{n_{\mathrm{arr}}=
 \frac{S_{yS}}{\sigma_{\mathrm{eq,arr}}}
 =\frac{130}{47{,}25}=2{,}75}
\]
\(e_b\) debe representar el borde susceptible de arrancarse
en la dirección de la fuerza. Recortes o cargas oblicuas cambian el modelo.
\cita{Mathcad: \mc{edge}, \mc{hole8}, \mc{TearG}, \mc{nTearG}.}
\end{frame}

\section{10. Raíz de cartela}
\gdef\mcref{Mathcad / sección 10}
\begin{frame}[t,label=mc10]{10. Flexión de la raíz de cartela}
Para una sección rectangular idealizada, con \(h_r=80\mm\):
\[
 Z_r=\frac{t_gh_r^2}{6},\qquad
 V_c=\sqrt{V_y^2+V_z^2}
\]
\begin{align*}
 \sigma_r&=\frac{M_c}{Z_r}
          =\frac{6M_c}{t_gh_r^2}=14{,}59\MPa,\\
 \tau_r&=\frac{V_c}{t_gh_r}=1{,}70\MPa.
\end{align*}
\[
 \sigma_{\mathrm{eq},r}
 =2\sqrt{\sigma_r^2+3\tau_r^2}=29{,}76\MPa,
 \qquad
 \boxed{n_r=4{,}37}
\]
\footnotesize
El factor 2 es provisional. El pliegue, el recorte y las fuerzas fuera del
plano exigen análisis local. También faltan precarga, presión bajo arandela
y ajuste de los casquillos.
\cita{[10] Mathcad: \mc{SigRoot}, \mc{TauRoot}, \mc{VMroot}, \mc{nRoot}.}
\end{frame}

\section{11. Cargas del piso}
\gdef\mcref{Mathcad / sección 11}
\begin{frame}[t,label=mc11]{11. Hipótesis de carga del piso}
\[
 W=300\mm,\quad D=125\mm,\quad t_f=2{,}5\mm
\]
\[
 \ell=65\mm,\quad a=40\mm,\quad
 P=100\N,\quad w_p=50\mm
\]
\[
 \boxed{q=\frac{2F_v}{WD}=0{,}03924\,\mathrm{N/mm^2}}
\]
Las líneas de apoyo se toman en los ejes de los refuerzos.
\(\ell\) es su separación y \(a\) es el voladizo delantero.

\textbf{Alcance de la presión elegida}
\[
 qWD+P=1\,571{,}5\N
 \qquad \text{frente a}\qquad F_v=735{,}75\N
\]
El factor 2 y el parche forman una envolvente local.
Su distribución debe justificarse con la bolsa llena y el contacto real.
\cita{Mathcad: \mc{Width}, \mc{Depth}, \mc{tf}, \mc{spanF}, \mc{over},
\mc{Ppatch}, \mc{strip}, \mc{qfloor}.}
\end{frame}

\begin{frame}[t]{11. Dos posiciones alternativas del parche}
Se analiza una franja de ancho unitario.
\begin{columns}[T,onlytextwidth]
\column{.48\textwidth}
\textbf{Tramo entre apoyos}
\[
 M_{\mathrm{tramo}}=
 \frac{q\ell^2}{8}+\frac{P\ell}{4w_p}
\]
\[
 M_{\mathrm{tramo}}=53{,}2236\N
\]
Carga uniforme y punto equivalente en el centro de una viga apoyada.
\column{.48\textwidth}
\textbf{Voladizo delantero}
\[
 M_{\mathrm{borde}}=
 \frac{qa^2}{2}+\frac{Pa}{w_p}
\]
\[
 M_{\mathrm{borde}}=111{,}392\N
\]
Carga uniforme y parche en el extremo libre.
\end{columns}
\bigskip
\[
 M_f=\max(M_{\mathrm{tramo}},M_{\mathrm{borde}}).
 \qquad [M_f]=\frac{\mathrm{N\,mm}}{\mathrm{mm}}=\mathrm{N}
\]
Se revisan posiciones alternativas de un parche.
\(M_f\) es momento por unidad de ancho.
\cita{Mathcad: \mc{Mspan}, \mc{Medge}, \mc{Mfloor}.}
\end{frame}

\section{12. Piso: tensión y espesor}
\gdef\mcref{Mathcad / sección 12}
\begin{frame}[t,label=mc12]{12. Desarrollo de la tensión del piso}
Para una franja de ancho \(u\):
\[
 I=\frac{ut_f^3}{12},\qquad
 c=\frac{t_f}{2},\qquad
 Z=\frac{I}{c}=\frac{ut_f^2}{6}
\]
El momento de la franja es \(M_fu\), por lo que:
\[
 \sigma_f=\frac{M_fu}{ut_f^2/6}
          =\boxed{\frac{6M_f}{t_f^2}}
\]
\[
 \sigma_f=\frac{6(111{,}392)}{2{,}5^2}
         =106{,}93632\MPa
\]
\[
 \boxed{n_f=\frac{130}{106{,}93632}=1{,}216<2{,}2}
\]
\alert{El piso actual queda por debajo del objetivo en este escenario.}
\cita{Mathcad: \mc{SigFloor}, \mc{nFloor}. El resultado no predice una rotura inmediata.}
\end{frame}

\begin{frame}[t]{12. Despeje del espesor orientativo}
\[
 \frac{S_{yS}}{6M_f/t_f^2}\ge n_{\mathrm{req}}
 \quad\Longrightarrow\quad
 t_f^2\ge\frac{6M_fn_{\mathrm{req}}}{S_{yS}}
\]
\[
 \boxed{t_{\mathrm{necesario}}=
 \sqrt{\frac{6M_fn_{\mathrm{req}}}{S_{yS}}}
 =3{,}363\mm}
\]
Comparación con \(t_{\mathrm{prop}}=3{,}5\mm\):
\[
 \sigma_{\mathrm{prop}}=
 \frac{6(111{,}392)}{3{,}5^2}=54{,}56\MPa,
 \qquad
 \boxed{n_{\mathrm{prop}}=2{,}383}
\]
También puede reducirse el voladizo o acercar un refuerzo al borde.

El CAD conserva \(2{,}5\mm\).
La comparación de \(3{,}5\mm\) pertenece a la memoria de cálculo.
\cita{Mathcad: \mc{tNeed}, \mc{tfProp}, \mc{SigProp}, \mc{nProp}.
Faltan contacto real, tolerancias y revisión local de la placa.}
\end{frame}

\section{13. Refuerzos y hombro}
\gdef\mcref{Mathcad / sección 13}
\begin{frame}[t,label=mc13]{13. Flexión del refuerzo inferior}
Tubo de \(20\times20\times1{,}5\mm\), con radios incluidos:
\[
 I_{\mathrm{ref}}=
 \operatorname{Ir}(b_r,b_r,2t_r)
 -\operatorname{Ir}(b_r-2t_r,b_r-2t_r,t_r)
\]
\[
 I_{\mathrm{ref}}=5\,827{,}22\mm^4,\qquad
 Z_{\mathrm{ref}}=\frac{I_{\mathrm{ref}}}{b_r/2}
                 =582{,}72\mm^3
\]
La viga apoyada recibe una fuerza equivalente \(\eta F_v\):
\[
 L_r=s+b=250\mm,\qquad M_{\max}=\frac{\eta F_vL_r}{4}
\]
\[
 \sigma_{\mathrm{ref}}
 =\frac{\eta F_vL_r}{4Z_{\mathrm{ref}}}
 =55{,}24\MPa,\qquad
 \boxed{n_{\mathrm{ref}}=4{,}34}
\]
\cita{Mathcad: \mc{Irib}, \mc{Zrib}, \mc{Lrib}, \mc{SigRib}, \mc{nRib}.
La comprobación no incluye las uniones en los extremos.}
\end{frame}

\begin{frame}[t]{13. Puente del anclaje de hombro}
El ligamento junto a la ranura se aproxima como una viga apoyada
con fuerza fuera del plano.
\[
 L_s=52\mm,\qquad b_s=24\mm,\qquad t_a=5\mm
\]
\[
 M_s=\frac{\eta F_vL_s}{4},
 \qquad
 Z_s=\frac{b_st_a^2}{6}
\]
\[
 \boxed{\sigma_s=
 K_{\mathrm{slot}}\frac{6(\eta F_vL_s/4)}{b_st_a^2}
 =100{,}43\MPa}
\]
\[
 \boxed{n_s=\frac{240}{100{,}43}=2{,}39}
\]
El espesor aparece al cuadrado porque la flexión ocurre fuera del plano.
Faltan cinta, costura, desgaste y radio de contacto en la ranura.
\cita{Mathcad: \mc{slot}, \mc{bridgeS}, \mc{ta}, \mc{SigSlot}, \mc{nSlot}.}
\end{frame}

\section{14. Cinturón y diagonales}
\gdef\mcref{Mathcad / sección 14}
\begin{frame}[t,label=mc14]{14. Puente del cinturón y raíz del yugo}
\begin{columns}[T,onlytextwidth]
\column{.48\textwidth}
\textbf{Puente de la ranura}
\[
 b_b=26\mm,\qquad t_b=5\mm
\]
\[
 \sigma_b=
 K_{\mathrm{slot}}
 \frac{6(\eta F_vL_s/4)}{b_bt_b^2}
\]
\[
 \sigma_b=92{,}70\MPa,\qquad
 \boxed{n_b=2{,}59}
\]
\column{.48\textwidth}
\textbf{Raíz del yugo}
\[
 Z_b=\frac{(100\mm)t_b^2}{6}
\]
\[
 \sigma_y=
 \frac{2\,\eta F_v(40\mm)}{Z_b}
 =98{,}88\MPa
\]
\[
 \boxed{n_y=2{,}43}
\]
\end{columns}
\medskip
El brazo de \(40\mm\) es una idealización local.
Quedan pendientes el equilibrio del arnés, la transferencia a la cadera
y los herrajes de liberación.
\cita{Mathcad: \mc{SigBelt}, \mc{nBelt}, \mc{Zbelt}, \mc{SigYoke}, \mc{nYoke}.}
\end{frame}

\begin{frame}[t]{14. Tracción de la diagonal activa}
\[
 h_X=H-30\mm-(120\mm+30\mm)=320\mm
\]
\[
 L_X=\sqrt{s^2+h_X^2}=377{,}36\mm
\]
Si \(\theta\) se mide desde la vertical:
\[
 \sin\theta=\frac{s}{L_X},\qquad
 F_X\sin\theta=F_h
 \quad\Longrightarrow\quad
 \boxed{F_X=F_h\frac{L_X}{s}=462{,}74\N}
\]
\[
 A_{\mathrm{net}}=(2\mm)(24-6{,}6)\mm=34{,}8\mm^2
\]
\[
 \sigma_X=\frac{F_X}{A_{\mathrm{net}}},
 \qquad
 \boxed{n_X=\frac{130}{\sigma_X}=9{,}78}
\]
Se acredita una diagonal traccionada. La opuesta puede pandear o aflojarse.
\cita{Mathcad: \mc{riseX}, \mc{LenX}, \mc{ForceX}, \mc{NetX}, \mc{nX}.
La cota de 120 mm requiere revisión explícita si cambia la altura de cuna.}
\end{frame}

\section{15. M6 y alcance}
\gdef\mcref{Mathcad / sección 15}
\begin{frame}[t,label=mc15]{15. M6 y mínimo de los márgenes incluidos}
\[
 R_6=K_b\frac{\eta F_v}{2}=309{,}015\N,
 \qquad A_{s6}=20{,}1\mm^2
\]
\[
 \sigma_{\mathrm{VM6}}=\sqrt{3}\frac{R_6}{A_{s6}}
 =26{,}63\MPa
 \qquad
 \boxed{n_{\mathrm{M6}}=\frac{450}{26{,}63}=16{,}90}
\]
La hoja reúne los márgenes de los modelos incluidos:
\[
 \mc{MinYield}=
 \min(n_R,n_{\mathrm{M8}},n_{\mathrm{bear}},n_f,\ldots)
 =1{,}216
\]
El mínimo ayuda a localizar el resultado más bajo.
Su alcance depende de los modos que se hayan incorporado.

Un margen alto del M6 deja pendientes la chapa, la unión completa
y la trayectoria de carga.
\cita{Mathcad: \mc{As6}, \mc{R6}, \mc{VM6}, \mc{nM6}, \mc{MinYield}.}
\end{frame}

\begin{frame}[t]{15. Comprobaciones calculadas y pendientes}
\begin{tabular}{@{}p{.25\textwidth}p{.68\textwidth}@{}}
\toprule
Fenómeno & Alcance de MC01\\
\midrule
Fluencia y flecha &
Modelos simplificados de riel, chapa, anclajes y tornillos.\\[3pt]
Pandeo &
Estimaciones elásticas ideales del riel. Faltan uniones e imperfecciones.\\[3pt]
Fatiga y fractura &
Sin vida calculada ni evaluación de grietas.\\[3pt]
Unión atornillada &
Faltan tracción, precarga, apertura, palanca y aflojamiento.\\[3pt]
Arnés y entorno &
Faltan textiles, costuras, ergonomía, temperatura y condiciones de uso.\\
\bottomrule
\end{tabular}
\bigskip
El margen mínimo de estas fórmulas describe las comprobaciones parciales.
La validación del conjunto requiere cerrar los mecanismos pendientes.
\end{frame}

\section{16. Resumen}
\gdef\mcref{Mathcad / sección 16}
\begin{frame}[t,label=mc16]{16. Resultados con las entradas iniciales}
\begin{columns}[T,onlytextwidth]
\column{.48\textwidth}
\small
\begin{tabular}{@{}lr@{}}
\toprule
Comprobación & Margen \(n\)\\
\midrule
Riel, \mc{nR} & 2,25\\
M8, \mc{nBolt8} & 5,31\\
Tubo, \mc{nBearR} & 2,68\\
Cartela, \mc{nBearG} & 2,32\\
Arrancamiento, \mc{nTearG} & 2,75\\
Raíz cartela, \mc{nRoot} & 4,37\\
\alert{Piso actual, \mc{nFloor}} & \alert{1,22}\\
\bottomrule
\end{tabular}
\column{.48\textwidth}
\small
\begin{tabular}{@{}lr@{}}
\toprule
Comprobación & Margen \(n\)\\
\midrule
Refuerzo, \mc{nRib} & 4,34\\
Hombro, \mc{nSlot} & 2,39\\
Puente cinturón, \mc{nBelt} & 2,59\\
Raíz yugo, \mc{nYoke} & 2,43\\
Diagonal, \mc{nX} & 9,78\\
M6, \mc{nM6} & 16,90\\
\bottomrule
\end{tabular}
\end{columns}
\bigskip
\[
 n_{\mathrm{req}}=2{,}2,
 \qquad
 \delta_R=3{,}440\mm,\quad\delta_{\mathrm{lim}}=5\mm
\]
Prioridad: revisar piso y apoyo delantero, y confirmar las zonas locales
del riel, cartelas y anclajes.
\end{frame}

\section{17. Actualización}
\gdef\mcref{Mathcad / sección 17}
\begin{frame}[t,label=mc17]{17. Medidas reales y actualización del modelo}
\begin{columns}[T,onlytextwidth]
\column{.49\textwidth}
\textbf{Datos por obtener}
\begin{itemize}
 \item Masa llena y componentes soportados.
 \item Forma de la bolsa llena y centro de carga.
 \item Apoyos y posiciones del arnés.
 \item Huella de contacto sobre el piso.
 \item Tubos, chapas y fijaciones reales.
\end{itemize}
\column{.47\textwidth}
\textbf{Revisión coordinada}
\begin{itemize}
 \item Guardar una copia con el cambio registrado.
 \item Ajustar entradas y brazos constantes.
 \item Recalcular y revisar unidades.
 \item Actualizar el CAD y comparar medidas.
 \item Completar detalles y pruebas de banco.
\end{itemize}
\end{columns}
\medskip
El CAD de referencia registra 54 diseños, 70 ocurrencias y aproximadamente
\(3{,}206\kg\), sin arnés completo.
La meta preliminar de \(1{,}5\) a \(2\kg\) aún requiere trabajo.

Los cambios en Mathcad y en Inventor se coordinan manualmente.
\end{frame}

\appendix
\section{Ampliaciones}
\gdef\mcref{Ampliación / fatiga}
\begin{frame}[t]{Fatiga: tensión media y amplitud}
Para un ciclo de tensión:
\[
 \boxed{\sigma_m=\frac{\sigma_{\max}+\sigma_{\min}}{2}},
 \qquad
 \boxed{\sigma_a=\frac{\sigma_{\max}-\sigma_{\min}}{2}}
\]
\begin{columns}[T,onlytextwidth]
\column{.48\textwidth}
\textbf{Información necesaria}
\begin{itemize}
 \item Historia de cargas al caminar.
 \item Número y secuencia de ciclos.
 \item Curvas S--N del detalle.
 \item Acabado, radios, agujeros y entorno.
\end{itemize}
\column{.48\textwidth}
\textbf{Alcance actual}

\(S_y\) permite revisar fluencia.
La vida en ciclos necesita datos cíclicos.

El \(K_t=2\) provisional de MC01 necesita revisión antes de emplearlo
en un análisis de fatiga.
\end{columns}
\cita{[11] MIT, Defect-Free Fatigue.
Esta ampliación explica un cálculo pendiente. MC01 no asigna vida útil.}
\end{frame}

\gdef\mcref{Ampliación / fractura}
\begin{frame}[t]{Fractura: presencia de una grieta}
Una aproximación lineal elástica de modo I utiliza:
\[
 \boxed{K_I=Y\,\sigma\sqrt{\pi a}}
\]
\begin{tabular}{@{}ll@{}}
\(a\) & Tamaño de grieta según la geometría elegida.\\[4pt]
\(Y\) & Factor geométrico.\\[4pt]
\(\sigma\) & Tensión nominal de referencia.\\[4pt]
\(K_I\) & Intensidad del campo de tensiones próximo a la grieta.\\
\end{tabular}
\bigskip

Para evaluar propagación inestable se necesita una tenacidad apropiada,
junto con espesor, condiciones de restricción y validez del modelo.

La memoria actual carece de esos datos. Esta expresión orienta
una ampliación futura cuando se defina el detalle fabricado.
\cita{[12] MIT, David Roylance, Introduction to Fracture Mechanics.}
\end{frame}

\gdef\mcref{Ampliación / sensibilidad}
\begin{frame}[t]{Sensibilidad del piso al espesor}
\[
 \sigma_f=\frac{6}{t_f^2}
 \left(\frac{qa^2}{2}+\frac{Pa}{w_p}\right)
\]
\begin{columns}[T,onlytextwidth]
\column{.49\textwidth}
\begin{tabular}{@{}rrr@{}}
\toprule
\(t_f\) [mm] & \(\sigma_f\) [MPa] & \(n_f\)\\
\midrule
2,5 & 106,94 & 1,22\\
3,0 & 74,26 & 1,75\\
3,5 & 54,56 & 2,38\\
4,0 & 41,77 & 3,11\\
\bottomrule
\end{tabular}
\column{.47\textwidth}
Con cargas fijas:
\[
 \sigma_f\propto\frac1{t_f^2}
\]
La presión contribuye con \(a^2\).
El parche contribuye con \(a\).

Acercar el apoyo reduce ambas contribuciones.
\end{columns}
\medskip
Pasar una placa plana de \(300\times125\mm\) de \(2{,}5\) a \(3{,}5\mm\)
agregaría aproximadamente \(101\) g, antes de considerar recortes o pliegues.
\end{frame}

\gdef\mcref{Consulta / nombres de Mathcad}
\begin{frame}[t]{Diccionario: cargas y rieles}
\begin{tabular}{@{}ll@{}}
\toprule
Símbolo de esta guía & Nombre en Mathcad\\
\midrule
\(m,g,a_v,a_h\) & \mc{Mass}, \mc{grav}, \mc{av}, \mc{ah}\\[3pt]
\(s,e,L\) & \mc{Sep}, \mc{ecc}, \mc{Le}\\[3pt]
\(A,A_m\) & \mc{Area}, \mc{Am}\\[3pt]
\(N_R,M_e,T_R\) & \mc{Nr}, \mc{Me}, \mc{Tr}\\[3pt]
\(\sigma_{\mathrm{ax}},\sigma_e,\sigma_h\) & \mc{SigAx}, \mc{SigEc}, \mc{SigH}\\[3pt]
\(\sigma_N,\sigma_{\mathrm{VM}}\) & \mc{SigN}, \mc{SigVM}\\[3pt]
\(\tau_T,\tau_V\) & \mc{TauT}, \mc{TauV}\\
\bottomrule
\end{tabular}
\bigskip

Unidades habituales: N, mm, MPa, mm\(^{2}\), mm\(^{3}\), mm\(^{4}\) y kg.
Los márgenes \(n\), los factores \(K\) y \(\nu\) son adimensionales.
\end{frame}

\begin{frame}[t]{Diccionario: uniones, piso y anclajes}
\begin{tabular}{@{}ll@{}}
\toprule
Símbolo de esta guía & Nombre en Mathcad\\
\midrule
\(M_c,p,a_H\) & \mc{MomC}, \mc{pitch}, \mc{armH}\\[4pt]
\(W,D,t_f\) & \mc{Width}, \mc{Depth}, \mc{tf}\\[4pt]
\(\ell,a,P,w_p\) & \mc{spanF}, \mc{over}, \mc{Ppatch}, \mc{strip}\\[4pt]
\(q,M_f,\sigma_f\) & \mc{qfloor}, \mc{Mfloor}, \mc{SigFloor}\\[4pt]
\(L_s,b_s,b_b\) & \mc{slot}, \mc{bridgeS}, \mc{bridgeB}\\[4pt]
\(L_X,F_X,A_{\mathrm{net}}\) & \mc{LenX}, \mc{ForceX}, \mc{NetX}\\
\bottomrule
\end{tabular}
\bigskip

Las variables conservan sus nombres originales en la hoja.
Esta tabla permite relacionarlos con los símbolos usados en las ecuaciones.
\end{frame}

\gdef\mcref{Consulta / fuentes}
\begin{frame}[t]{Fuentes de Mathcad y materiales}
\footnotesize
\begin{description}
 \item[{[1]}] Ayuda local de Mathcad clásico, carpeta \mc{doc/HELP_ES/Help}.
 Operadores, funciones, unidades y recálculo.
 \href{https://www.ptc.com/en/support/help/mathcad}{Centro oficial de ayuda PTC}.
 \item[{[3]}] \href{https://www.hydro.com/globalassets/01-products--services/extruded-profiles/americas/ena-resources/alloy-data-sheets/hydro_2019_data_sheet_6061.pdf}{Hydro, Alloy 6061.}
 Propiedades del perfil según temple y producto.
 \item[{[4]}] \href{https://atlassteels.com.au/wp-content/uploads/2021/08/Aluminium-Alloy-5052-Data-Sheet-24-06-21.pdf}{Atlas Steels, Aluminium Alloy 5052 Data Sheet.}
 \item[{[5]}] \href{https://www.atlassteels.com.au/documents/TN14-5052\%20and\%205251\%20rev\%20June\%202013.pdf}{Atlas Steels, Technical Note 14, 5052 and 5251.}
 Diferencias entre especificaciones.
 \item[{[6]}] \href{https://www.bossard.com/-/media/bossard-group/website/documents/technical-resources-old/stainless-steel-fasteners.pdf}{Bossard, Stainless steel fasteners.}
 Propiedades de fijaciones según ISO 3506.
\end{description}
\medskip
Las ecuaciones y valores iniciales corresponden a MC01.
Las propiedades del material comprado deben sustituir los valores provisionales.
\end{frame}

\begin{frame}[t]{Fuentes de diseño mecánico}
\footnotesize
\begin{description}
 \item[{[2]}] \href{https://ocw.mit.edu/courses/2-080j-structural-mechanics-fall-2013/30dc1a0f74debf21fb92a1df56616929_MIT2_080JF13_Lecture12.pdf}{MIT, Structural Plasticity.} Von Mises y Tresca.
 \item[{[7]}] \href{https://ocw.mit.edu/courses/2-080j-structural-mechanics-fall-2013/3533c046dafcc488f1432e92d05ec210_MIT2_080JF13_Lecture5.pdf}{MIT, Solution Method for Beam Deflections.}
 \item[{[8]}] \href{https://ocw.mit.edu/courses/16-001-unified-engineering-materials-and-structures-fall-2021/mit16_001_f21_lec33lec34lec35.pdf}{MIT, Buckling of Beams.}
 \item[{[9]}] \href{https://ocw.mit.edu/courses/2-080j-structural-mechanics-fall-2013/345c462fe40421ea0f795bd35d2e4af1_MIT2_080JF13_Lecture11.pdf}{MIT, Buckling of Plates and Sections.}
 \item[{[10]}] \href{https://ntrs.nasa.gov/api/citations/19900009424/downloads/19900009424.pdf}{NASA RP-1228, Fastener Design Manual, 1990.}
 Referencia general para uniones.
 \item[{[11]}] \href{https://ocw.mit.edu/courses/2-002-mechanics-and-materials-ii-spring-2004/db45866e2fc0cf164d96828936748e8f_lec21_notes.pdf}{MIT, Anand y Parks, Defect-Free Fatigue.}
 \item[{[12]}] \href{https://web.mit.edu/course/3/3.11/www/modules/frac.pdf}{MIT, David Roylance, Introduction to Fracture Mechanics.}
\end{description}
\medskip
\textbf{Documento de referencia:} \mc{Memoria_SEDAC_P2_02_MC01_EDITABLE.xmcd}.

\textbf{Estado:} prediseño P2.02. Las fórmulas parciales se revisarán
con mediciones, detalles fabricables y pruebas de banco.
\end{frame}
\end{document}
